\section{Iterarea prin submăștile unei măști}

Ocazional avem nevoie ca, pentru o mască binară $m$, să iterăm prin toate numerele binare care conțin biți 1 doar pe poziții pe care $m$ conține biți 1. De exemplu, pentru masca \texttt{0110010}, dorim să iterăm prin măștile

\begin{verbatim}
0000000
0000010
0010000
0010010
0100000
0100010
0110000
0110010 = M
\end{verbatim}

Surprinzător, putem face acest lucru cu o buclă elementară:

\begin{minted}{c}
for (unsigned s = m; s; s = (s - 1) & m) {
  // s este o submască a lui m
}
\end{minted}

Remarcăm că bucla nu tratează masca \ccode{s = 0}. Dacă o doriți și pe aceea, faceți un caz particular în afara buclei sau scrieți un cod puțin mai urît pentru a include 0 în buclă, de exemplu:

\begin{minted}{c}
int s = m + 1;
do {
  s = (s - 1) & m;
  // procesează masca s
} while (s);
\end{minted}

De ce funcționează acest cod? În cazul simplu cînd $m$ este contiguă, să zicem \texttt{1111111}, observăm că \ccode{(s - 1) & m = s - 1}, deci \ccode{s} iterează pur și simplu descrescător de la \ccode{1111111} la \ccode{0000000}. Dacă $m$ conține și zerouri intercalate, atunci $s$ va itera descrescător prin combinațiile de 1, ignorînd zerourile. Exemplu:

\begin{verbatim}
m           = 110111110011000
s           = 010110000000000
s - 1       = 010101111111111
(s - 1) & m = 010101110011000
\end{verbatim}
